\section{Discussion}

\subsection{Space Matters More Than Abstract Disinterest}

The clearest finding of this study is that the female college students in the sample were not universally inactive, nor did they completely reject the gym; rather, they exhibited a pronounced threshold for entry specifically into the free-weight area.
Although 78.13\% of participants had used the campus gym and 46.88\% exercised 2--3 times per week, only 15.63\% most frequently used the free-weight area.
This discrepancy suggests that the issue should not be interpreted as ``women lack exercise interest,'' but rather as uneven participation resulting from spatial differentiation within the campus gym.

The mean score for spatial pressure/training anxiety reached 3.72, with discomfort when many men were present, pressure from crowded conditions, fear of incorrect form, and perceptions of male dominance all at moderate-to-high levels.
After supplementing the participant-level data, spatial pressure/training anxiety showed a strong positive correlation with the frequency of actively avoiding the free-weight area ($\rho = .65$), and regular users of the free-weight area also had notably lower spatial pressure scores than non-regular users.
Although this bivariate relationship cannot demonstrate causal direction, it reinforces the core explanation of this paper: avoidance is not an isolated attitude but co-occurs with specific spatial pressures.
This result resonates with \textcite{turnock2021gyms}'s finding that women in mixed-gender gyms feel tolerated rather than genuinely welcomed, and aligns with \textcite{rendall2024gym}'s identification of gym intimidation as an exercise barrier for female college students.
From the spatial theory perspective of \textcite{lefebvre1991production}, the free-weight area is not a neutral equipment zone; rather, it is produced---through mirrors, open visibility, bodily display, weight disparities, and gender ratios---as a social space with an entry threshold.

The interviews further revealed that this threshold does not arise solely from hostility on the part of individual men.
On one hand, interviewees acknowledged that most people were ``just doing their own workout''; on the other hand, they still felt overly visible due to mirrors, peripheral vision, equipment occupation, and unfamiliarity with exercises.
This aligns with the self-surveillance mechanism emphasized by objectification theory: even in the absence of explicit external evaluation, women may preemptively position themselves as being watched in visible spaces \parencite{fredrickson1997objectification,clark2018gaze}.
Thus, avoiding the free-weight area is better understood as a form of anticipatory management of visibility and the possibility of error, rather than simple timidity.
Notably, training self-efficacy showed virtually no correlation with avoidance frequency ($\rho = -.11$, $p = .562$), despite a moderate self-efficacy mean ($M = 3.49$); this suggests that the barrier lies not in participants' general confidence about achieving training goals but in the specific spatial and beginner-entry conditions of the free-weight area---further supporting the interpretation that avoidance is situationally driven rather than attributable to individual dispositions.

\subsection{Beginner Access as the Primary Barrier, With Indirect Media Effects}

The initial research framework hypothesized that social media appearance internalization might significantly drive resistance training avoidance, particularly through the concern that strength training would cause the body to deviate from the slender feminine ideal.
The actual data did not fully support this.
Although 59.38\% of participants agreed in Q21 that appearance content influenced their exercise choice, the mean for Q22 (``worry that strength training will not look slender enough'') was only 1.85, with 81.25\% in the lower score range.
Among avoidance reasons, ``do not want to build visible muscle'' also registered at 0\%.
The individual-level correlation similarly showed only a weak positive trend between social media appearance internalization and avoidance frequency ($\rho = .31$, $p = .084$), which was less robust than the spatial pressure variables.

This does not mean that social media is unimportant; rather, it indicates that its role in this sample was more indirect and complex.
Q21's mean reached 3.40 with an upper score range of 59.38\%, indicating that appearance content does influence the exercise selection process, but Q22's low scores suggest that this influence operates more at the goal level (e.g., making fat loss the most common exercise goal) without directly translating into a strong fear of muscular outcomes from strength training.
Both P2 and P3 in the interviews mentioned encountering related debates but did not agree with the claim that ``women who do strength training look bad.''
This is consistent with changes revealed by recent research on fitspiration and the fit-ideal: contemporary female body ideals are not exclusively about thinness but simultaneously incorporate elements of health, toning, strength, and displayability \parencite{tiggemann2015fitspiration,donovan2020fitideal}.

By contrast, the questionnaire and interviews more consistently pointed to the problem of ``how beginners enter.''
``Do not know how to use equipment'' and ``too many men'' were equally the most frequent avoidance reasons, and ``do not know where to start'' also reached 50.00\%.
Participants with systematic strength training experience had lower avoidance frequency than those without systematic experience; although limited by sample size, this result is consistent with the interview narrative that ``having learned/having someone teach'' can reduce uncertainty.
P2 began trying strength training after learning through a course, while P3, not knowing the equipment and movements, primarily performed bodyweight exercises in her dormitory.
This suggests that if the university only provides gym equipment without offering entry paths that beginners can understand, practice, and make mistakes in, formal resource accessibility will not automatically translate into actual participation.

\subsection{Strategic Negotiation Rather Than Passive Withdrawal}

This study also suggests that participants should not be portrayed as passive subjects entirely suppressed by spatial and appearance norms.
Many avoidance behaviors were actually strategic: going during off-peak hours, going when fewer people were present, finding a friend to go with, switching to the cardio or fixed machine area, and observing the environment before deciding whether to continue.
These practices show that participants were not uninterested in exercise but were seeking relatively controllable ways to enter a space where they did not feel entirely at ease.

This ``negotiated participation'' echoes \textcite{xiong2021fanshen}'s discussion of Chinese women's fitness narratives: women's fitness cannot be reduced to a singular narrative of ``being disciplined'' or ``being empowered,'' but involves continuous negotiation among body image, health management, social relationships, and self-identity.
In this study, women's participation in fitness could be organized by appearance norms, but it could also become a practice for enhancing physical capability, building self-care, and asserting spatial presence.
Specifically, P2 understood strength training as health and posture improvement, and P3, despite being nervous, still said she wanted to go to the gym in the future.
Therefore, resistance training avoidance should not be understood as a fixed identity but as a behavioral state that is repeatedly adjusted under specific spatial conditions.

It is worth noting that RQ3 originally focused on the negotiation between ``health motivation'' and ``digital appearance norms,'' but the data from this sample more consistently pointed to a different kind of negotiation---participants were primarily not choosing between health and appearance but rather seeking strategic ways to participate in a space that made them feel overly visible and lacked entry paths.
This shift itself corroborates the earlier findings: in this sample, spatial visibility and beginner barriers constituted more prominent obstacles than appearance pressure, and the core target of ``negotiation'' was spatial conditions rather than appearance norms.

\subsection{Practical Implications}

In practical terms, the most direct recommendation is not simply to encourage female students to ``be braver,'' but to lower the entry costs for beginners in the free-weight area.
The most popular intervention in the questionnaire was the beginner-friendly women's strength training workshop, with an upper score range reaching 96.88\%.
The correlation between spatial pressure/training anxiety and the intervention preference index was also high ($\rho = .67$), indicating that the demand for supportive interventions primarily came from participants who already perceived the free-weight area as difficult to enter.
Such workshops should focus on basic movements, equipment usage, weight selection, training programming, and common errors, rather than merely promoting body sculpting.
If these workshops could be led by female coaches, student assistants, or trained peer mentors, they might further reduce beginner discomfort in male-dominated spaces.

Second, spatial design should reduce the pressure of ``being seen as soon as you enter.''
Support for semi-private partitioning or visual screening was high (upper score range 81.25\%), and the interviews explicitly mentioned partitions, small spaces, fewer large mirrors, and greater attention to cleanliness.
This does not necessarily require complete gender segregation of the gym; rather, it can be achieved by distributing lighter equipment, setting up beginner practice corners, improving the relationship between mirrors and traffic flow, and providing clear equipment instructional videos, thereby creating lower-risk practice areas for beginners.

Third, campus communication should shift from ``slimming/fat loss'' toward functional, body-diverse, and beginner-friendly narratives.
Although this sample did not strongly worry that strength training would make them not slender enough, fat loss/slimming remained the most common exercise goal.
If university communications continue to present exercise as an appearance management project, they may reinforce body comparison; a more appropriate approach would be to emphasize the benefits of strength training for physical fitness, posture, bone health, menstrual health, and daily functional capacity, while showcasing women at different training levels and body types.

\subsection{Theoretical Reflection}

The three-layer integrated theoretical framework proposed in the literature review---spatial production, body objectification, and behavioral outcomes---effectively organized the findings of this study overall, but the level of support for each layer was asymmetric.
The first layer (spatial production) received the strongest support: the strong correlation between spatial pressure/training anxiety and avoidance frequency ($\rho = .65$), together with the interview narratives about mirrors, visibility, and the male-dominated atmosphere, jointly corroborated that the free-weight area is produced---through social relations---as a space with an entry threshold.
The second layer (appearance internalization) played a weaker role than initially anticipated: the correlation between social media appearance internalization and avoidance did not reach significance, and ``fear of getting bulky'' was virtually absent as an avoidance reason; however, Q21's high scores suggested that appearance content still influenced exercise goal selection, albeit operating primarily at the goal level rather than directly driving avoidance behavior.
The third layer (behavioral outcomes) was further enriched by the data: participants' behavior manifested not merely as a binary choice between avoidance and participation but as strategic negotiation situated between the two, which is consistent with \textcite{xiong2021fanshen}'s observation of ``neither fully disciplined nor fully empowered'' in women's fitness narratives.
This asymmetry is itself a finding---it suggests that among female college students lacking systematic strength training experience, spatial visibility and beginner barriers may be more proximal and direct obstacles than appearance internalization.
The mixed-methods design strengthened this interpretation through triangulation: the quantitative correlation ($\rho = .65$) between spatial pressure and avoidance established the direction and strength of the association, while the interview narratives provided the contextual mechanism---mirrors, male numerical dominance, fear of incorrect form, and equipment unfamiliarity---explaining why spatial pressure translates into avoidance behavior.
Conversely, the non-significant correlation between social media appearance internalization and avoidance ($\rho = .31$, $p = .084$) was enriched by the qualitative finding that interviewees' relationships with media content were differentiated and often dismissive, lending depth to what would otherwise be a merely negative statistical result.

\subsection{Limitations}

The greatest limitation of this study is the small sample size and single-campus context: the core questionnaire sample consisted of only 32 participants, and interviews involved only 3 individuals.
After supplementing the participant-level response matrix, this paper was able to calculate internal consistency, individual-level construct scores, Spearman correlations, and exploratory group comparisons; however, these results remain susceptible to the influence of outliers, group imbalance, and multiple comparisons.
Therefore, this paper does not conduct regression, mediation, or predictive modeling, nor does it interpret $p$ values as strong evidence.

Second, in some multiple-choice items, a minority of options were blank in the aggregate export table; this paper only reports options with explicit frequencies or those inferable from the valid sample size in the descriptive frequencies, and retains missing markers in the analysis script.
Third, the scale items were context-specific self-developed or adapted short items; although they drew on the theoretical dimensions of SEAM and SATAQ-4, they have not undergone full psychometric validation.
In particular, the three intervention preference items measured different policy/spatial proposals and should not be over-interpreted as a single psychological scale.
Fourth, this study used cross-sectional self-report data, which cannot determine whether spatial pressure causes avoidance or whether students who already avoid the free-weight area are more likely to evaluate that space as intimidating in their imagination.

Future research should preserve participant-level raw response data, expand the sample to multiple universities, and employ fully validated scales or formally validated Chinese short-form instruments.
If conditions permit, researchers could compare free-weight area usage before and after participation in a beginner-friendly women's strength training workshop, thereby more directly evaluating whether workshops and spatial interventions can reduce resistance training avoidance.
